\documentclass{article}
\usepackage[UTF8]{ctex}
\usepackage{amsmath,amssymb,amsfonts}
\usepackage{enumitem}
\usepackage{geometry}
\geometry{a4paper, margin=1in}
% === 错题本排版优化（不参与题目标识） ===
\setlength{\parindent}{0pt}
\setlength{\parskip}{0.35em plus 0.10em minus 0.08em}
\setlength{\emergencystretch}{3em}
\setlist{itemsep=0.25em, topsep=0.40em, parsep=0pt}
\allowdisplaybreaks[2]
% === 错题本排版优化结束 ===
\begin{document}
\section*{吉林大学2025年高等代数真题}
\section*{69．吉林大学2025年高等代数与解析几何试卷}

\subsection{题目1}
一．设 \(f(x), g(x)\) 都是数域 \(\Omega\) 上的首一多项式，且满足 \(\left(f^{2024}, g^{2025}\right)=\left(f^{2025}, g^{2024}\right)\) 。证明：对任意的正整数 \(m, n\) ，有 \(\left(f^{m}, g^{n}\right)=\left(f^{n}, g^{m}\right)\) 。

\subsection{题目2}
二．设 \(V\) 是数域 \(\Omega\) 上的 \(n\) 维向量空间，\(n \geq 2 . \sigma\) 是 \(V\) 上的线性变换．证明：存在 \(V\) 的线性变换 \(\tau\) ，使得 \(\operatorname{dim} \sigma(V)=\operatorname{dimKer} \tau\) ，且 \(\sigma \tau=\tau \sigma=0\) ．

\subsection{题目3}
三。设 \(A, B\) 都是 \(n\) 阶实对称矩阵，证明下列陈述（1）与（2）等价。

\begin{itemize}
  \item[（1）] 对任意正整数 \(m, k\) ，都有 \(m A+k B\) 是正定矩阵。
  \item[（2）] \(A, B\) 都是半正定矩阵，且 \(A X=0\) 与 \(B X=0\) 的解空间的交为零空间．
\end{itemize}

\subsection{题目4}
四．设 \(V\) 是 \(n\) 维欧氏空间，\(n \geq 2, \sigma\) 是 \(V\) 的对称变换，记\(L=\{\tau\) 是 \(V\) 的线性变换 \(\mid \sigma \tau=\tau \sigma\}\) 。\\
证明： \(\operatorname{dim} L=n\) 当且仅当 \(\sigma\) 的每个特征子空间都是一维的。

\subsection{题目5}
五．设 \(V\) 是 \(n\) 维复向量空，\(n \geq 2, \sigma, \tau\) 是 \(V\) 的线性变换，假设 \(\sigma\) 可对角化，\(\tau\) 是幂零变换，且 \(\sigma \tau=\tau \sigma\)。\\
（1）证明：\(\sigma\) 与 \(\sigma+\tau\) 的特征多项式相同．\\
（2）证明：\(\tau\) 与 \(\sigma+\tau\) 的极小多项式的次数相等当且仅当存在复数 \(a\) ，使得对任意的 \(u \in V\) ，有\(\sigma(u)=a u\) ．

\subsection{题目6}
六．设 \(a, \beta, \gamma, \delta\) 是几何空间中的四个向量，证明：

\[
(\alpha \times \beta) \cdot(\gamma \times \delta)-(\alpha \times \gamma) \cdot(\beta \times \delta)-(\delta \times \alpha) \cdot(\beta \times \gamma)=0 .
\]

\subsection{题目7}
七．在空间直角坐标系中，求二次曲面 \(x^{2}+2 x y-3 y^{2}=2 z\) 上经过点 \((1,-1,-2)\) 的所有直母线的标准方程．

\subsection{题目8}
八．在空间直角坐标系中，设单叶双曲面 \(S\) 的一个一般方程为 \(x^{2}+\frac{y^{2}}{4}-\frac{z^{2}}{9}=1\) 。证明：存在经过 \(y\) 轴的平面，使得该平面与单叶双曲面 \(S\) 的截线为＂长轴长：短轴长 \(=\sqrt[3]{3}\) 的椭圆＂，并求出所有满足上述条件的平面（写出平面的一般方程即可）。
\end{document}
